\begin{proof}
$\rho(w_1)\rho(w_2)+\rho(w_2)\rho(w_1)=2(w_1,w_2)$ applied to $\xi\ne0$ gives $(w_1,w_2)=0$. By Corollary \ref{cor:c7-Q}, for $(w_1,w_2)=0$ the component is a multiple of $\int\omega\,\big[(\iota_{w_1}A,(\rho(w_2)\xi,\lambda))+(\iota_{w_2}A,(\rho(w_1)\xi,\lambda))\big]$, which vanishes; the classical component is the same functional (proof of Corollary \ref{cor:c7-Q}). In the evaluation of Theorem \ref{thm:c11-bilinears} the two-body part on these bilinears, $-2L\,\mathcal B_1+4L(\mathcal B_2+\mathcal B_3)$ units, vanishes term by term, with $\mathcal B_1=\{S^{(1)}(\xi),O\}=0$ and $(w_1,w_2)=0$, and the one-body term is a multiple of $\mathcal B_1$; by Corollary \ref{cor:c11-F0} the same holds with the computed solution and every free direction, the coefficients of $\mathcal B_1,\dots,\mathcal B_4$ multiplying functionals that vanish here, and $\hbar L_1(O)^{(0)}$ vanishes with $(w_1,w_2)$. The component along $(\xi,\xi)$ of a functional of degree $0$ and Weyl weight $2$ has degree $-1$ and Weyl weight $3$ (Proposition \ref{prop:c7-select}, applied once for each argument), so an output with two letters contains an antifield, of Weyl weight at least $\tfrac52$, and a field of weight at least $1$, of total weight at least $\tfrac72>3$; the one-body term $-\{\hbar F_1^{(2)}(\xi,\xi),O\}$ vanishes, a quadratic functional of degree $-2$ and Weyl weight $5$ having the legs $(\lambda^*,\lambda^*)$ or $(A,c^*)$ and no leg $A^*$. $\Phi$ maps classes to classes, so a vanishing differential stays zero under every normalisation.
\end{proof}
